\section{Ethics and Limitations}
\label{sec:ethics-limitations}

All scenarios, entities, traces, policies, and tool outputs are project-authored
synthetic data, not customer records or private company policies; no real person
is intentionally represented. Nevertheless, future free-text annotations could
introduce personal information, so annotators use participant codes, may not paste
real credentials or records, and exports require PII/secret review. The tasks
simulate consequential purchases, refunds, identity and access checks,
accessibility requests, and travel-document workflows. Their simplified policies
are unsuitable for operational, legal, medical, financial, employment, or other
decisions about real people.

The sealed split is synthetic and may contain lexical regularities, simplified
jurisdictional assumptions, and limited real-world coverage. Human reconstruction
and deterministic checks establish protocol validity, not ecological validity.
The 30.5\% oracle-sequence-plus-enum-first diagnostic reveals ordered-enum
regularity under an evaluator-provided workflow path; it is not name-only
performance. Moreover, 69 normalized-topology overlaps with development tasks
preclude a claim of complete structural holdout, although there are no exact
development identity, trace, or action-graph collisions. A
post-seal replication can adapt separately licensed ToolSandbox scenarios and
the telecom subset of $\tau^2$-Bench \citep{lu2025toolsandbox,barres2025tau2},
but must report each source separately and may not replace rejected primary
families after outcomes are observed.
Moreover, the repository does not yet declare a top-level data license; no open
license should be inferred before release terms are set. The sealed tasks cannot
be repaired or filtered after model outcomes are observed. Accordingly, this
draft reports completed human validation and construction diagnostics, but no
confirmatory performance or representativeness claim.
